\documentclass[11pt,a4paper]{article}

% --- Variables du TP
\newcommand{\TPModule}{Datawarehouse}
\newcommand{\TPNumero}{1}
\newcommand{\TPKeywords}{data engineering, ingestion, dbt, bronze, silver}
\newcommand{\TPSujet}{Mise en place du pipeline décisionnel}
\newcommand{\TPTheme}{Clinique}
\usepackage[utf8]{inputenc}
\usepackage[french]{babel}
\usepackage[T1]{fontenc}
\usepackage{amsmath}
\usepackage{amsfonts}
\usepackage{xcolor}
\usepackage{amssymb}
\usepackage{graphicx}
\usepackage[left=2cm,right=2cm,top=2cm,bottom=2cm]{geometry}
\usepackage{fancyhdr}
\usepackage[bottom]{footmisc}
\usepackage{tikz}
\usepackage{fourier}
\usepackage{pst-text}
\usepackage{multicol}
\usepackage{comment}
\usepackage{array}
\usepackage{enumitem}
\setlist[itemize]{label=\textbullet} % pour ne plus utiliser les bulletpoints après \item
\setlength{\parindent}{0pt} % pour supprimer les \noindent
\usepackage{amssymb} % pour les symboles
\usepackage{tabularx}
\usepackage{tcolorbox}
\usepackage{fontawesome}
\usepackage{listings}

% --- Réglages listings (inchangé)
\lstset{
  basicstyle=\ttfamily\small,
  backgroundcolor=\color{gray!10},
  frame=single,
  rulecolor=\color{gray!40},
  breaklines=true,
  columns=fullflexible,
  showstringspaces=false,
  aboveskip=0.5em,
  belowskip=0.5em,
  keywordstyle=\color{blue},
  commentstyle=\color{gray},
}
\lstdefinelanguage{markdown}{
  sensitive=true,
  morecomment=[l]{<!--},
  morecomment=[s]{<!--}{-->},
  morekeywords={},
}

% --- Réglages listings (inchangé)
\lstset{
  basicstyle=\ttfamily\small,
  backgroundcolor=\color{gray!10},
  frame=single,
  rulecolor=\color{gray!40},
  breaklines=true,
  columns=fullflexible,
  showstringspaces=false,
  aboveskip=0.5em,
  belowskip=0.5em,
  keywordstyle=\color{blue},
  commentstyle=\color{gray},
}
\lstdefinelanguage{markdown}{
  sensitive=true,
  morecomment=[l]{<!--},
  morecomment=[s]{<!--}{-->},
  morekeywords={},
}

% --- Fancyhdr construit depuis les variables
\pagestyle{fancy}
\setlength{\headheight}{26pt}
\renewcommand{\baselinestretch}{1.3}

\fancyhead[L]{\textbf{EPSI\\\TPModule}}
\fancyhead[R]{\textbf{TP n°\TPNumero\\\TPSujet}}
\renewcommand{\arraystretch}{1}

% --- Hyperref + metadata construits depuis les variables
\usepackage[
  colorlinks=true,
  linkcolor=blue!70!black,
  urlcolor=blue!70!black,
  citecolor=blue!70!black
]{hyperref}

\hypersetup{
  pdftitle={TP n°\TPNumero\space - \TPSujet},
  pdfauthor={Thibault CLEMENT},
  pdfsubject={TP n°\TPNumero\space - \TPSujet},
  pdfkeywords={\TPKeywords},
}

% --- Composants réutilisables
\newcommand{\TPFrontPage}{%
  ~
  \begin{center}
    \framebox{\parbox[c]{11cm}{%
      \begin{center}
        \Large{\textbf{TP n°\TPNumero~:~\TPSujet\\\large{\TPTheme}}}
      \end{center}
    }}
  \end{center}
  \vspace{0.5cm}
}

% --- Tes macros custom existantes
\newcommand{\file}[1]{~\texttt{\textcolor{blue!70!black}{#1}}~}

\newtcolorbox{reflectionbox}[1][]{%
  colback=blue!5,
  colframe=blue!60,
  title={\faLightbulbO~~Réflexions},
  fonttitle=\bfseries,
  left=1mm,
  right=1mm,
  top=1mm,
  bottom=1mm,
  #1
}


\newcommand{\checkbox}{\(\square\)~~}
\graphicspath{{images/}}
\setlength{\columnseprule}{1pt}
\def\columnseprulecolor{\color{black}}



\begin{document}
\TPFrontPage

% ---------------------------------------------------------------------------------------------

\section*{Contexte métier}

Vous intervenez pour une clinique privée.

Le directeur souhaite mieux piloter l’activité de ses services :
\begin{itemize}
  \item tension dans certains services,
  \item refus de patients,
  \item impact sur la satisfaction,
  \item gestion du personnel,
  \item rentabilité des lits.
\end{itemize}

Les données existent, mais elles proviennent de plusieurs systèmes opérationnels
et ne sont pas encore intégrées dans une architecture décisionnelle.

Votre mission consiste à mettre en place les \textbf{fondations du pipeline de données}
afin de permettre, dans un second temps, la production d’indicateurs décisionnels.

\bigskip

\section*{Questions stratégiques du directeur (fil rouge)}

À terme, le directeur souhaite pouvoir répondre aux questions suivantes :

\begin{enumerate}
  \item Quels services sont le plus souvent en tension ?
  \item Quel est le taux moyen d’occupation par service ?
  \item Combien de patients sont refusés chaque semaine et pourquoi ?
  \item La satisfaction diminue-t-elle en période de forte activité ?
  \item Existe-t-il un lien entre présence du personnel et tension ?
  \item Quelle est la durée moyenne de séjour par service ?
  \item Quels services semblent les plus rentables ?
  \item Peut-on anticiper les semaines critiques ?
\end{enumerate}

\textbf{Important :}
Ce TP ne vise pas encore à répondre à ces questions.
Il vise à construire les \textbf{fondations techniques} permettant d’y répondre dans le TP2.

% ---------------------------------------------------------------------------------------------

\section*{Objectifs pédagogiques}

Ce TP a pour objectif de mettre en pratique :

\begin{itemize}
  \item le rôle des couches \textbf{RAW, BRONZE, SILVER},
  \item l’ingestion reproductible de données,
  \item la structuration d’un pipeline ELT,
  \item l’utilisation de \textbf{dbt} pour transformer les données,
  \item la mise en place de tests de qualité,
  \item le principe du chargement incrémental.
\end{itemize}

% ---------------------------------------------------------------------------------------------

\section*{Jeux de données fournis}

Les fichiers suivants sont disponibles dans le dossier \file{data/} :

\begin{itemize}
  \item \file{patients.csv}
  \item \file{staff.csv}
  \item \file{staff\_schedule.csv}
  \item \file{admissions\_requests.csv}
\end{itemize}

Ces fichiers constituent la \textbf{couche RAW}.

\bigskip

\textbf{Important :}
Aucune table agrégée hebdomadaire n’est fournie.
La construction d’indicateurs sera réalisée dans le TP2.

% ---------------------------------------------------------------------------------------------

\section{Étape 1 : Exploration et Data Profiling}

\subsection*{Objectif}

Comprendre la structure, le contenu et les spécificités des données sources
avant toute industrialisation.

\subsection*{Consignes}

Cette étape peut être réalisée dans un \textbf{notebook}.

\begin{enumerate}
  \item Charger les différents fichiers CSV.
  \item Analyser :
  \begin{itemize}
    \item types de colonnes,
    \item valeurs manquantes,
    \item formats de dates,
    \item cohérence des identifiants,
    \item distributions simples.
  \end{itemize}
  \item Identifier les relations potentielles entre les tables.
  \item Déterminer le grain naturel de chaque fichier.
  \item Construire un dictionnaire de données synthétique.
\end{enumerate}


% ---------------------------------------------------------------------------------------------

\section{Étape 2 : Mise en place de la couche BRONZE}

\subsection*{Objectif}

Mettre en place un stockage brut, traçable et reproductible des données sources.

\subsection{Réflexion}

\begin{enumerate}
  \item Expliquer le rôle de la couche BRONZE dans une architecture médaillon.
  \item Pourquoi ne faut-il pas appliquer de logique métier à ce stade ?
  \item Quelle différence entre RAW et BRONZE ?
\end{enumerate}

\subsection{Implémentation}

\begin{enumerate}
  \item Créer une base SQLite : \file{dwh/clinic\_dwh.db}.
  \item Écrire un script Python \file{scripts/ingest\_bronze.py} :
  \begin{itemize}
    \item lecture des fichiers CSV depuis \file{data/},
    \item chargement dans SQLite,
    \item une table par fichier.
  \end{itemize}
\end{enumerate}

\subsection*{Règles d’ingestion}

\begin{itemize}
  \item Préfixe des tables : \file{bronze\_}\\\textit{Note : SQLite ne gérant pas nativement les schémas (bronze, silver, gold), la séparation des couches sera réalisée par convention de nommage.}
  \item Aucune transformation métier.
  \item Ajout des colonnes techniques :
  \begin{itemize}
    \item \file{batch\_id}
    \item \file{ingestion\_ts}
    \item \file{source\_file\_name}
  \end{itemize}
\end{itemize}


% ---------------------------------------------------------------------------------------------

\section{Étape 3 : Transformation vers la couche SILVER (ELT avec dbt)}

\subsection*{Objectif}

Produire des tables propres, typées et cohérentes,
prêtes pour une modélisation décisionnelle.

\subsection{Réflexion ETL vs ELT}

\begin{enumerate}
  \item Identifier au moins trois transformations nécessaires.
  \item Décrire une approche ETL.
  \item Décrire une approche ELT.
  \item Comparer les deux approches :
  flexibilité, traçabilité, performance, maintenabilité.
\end{enumerate}

\subsection{Implémentation avec dbt}

\begin{enumerate}
  \item Configurer dbt pour utiliser \file{clinic\_dwh.db}.
  \item Créer un modèle par table BRONZE :
  \begin{itemize}
    \item \file{silver\_patients}
    \item \file{silver\_staff}
    \item \file{silver\_staff\_schedule}
    \item \file{silver\_admissions\_requests}
  \end{itemize}
  \item Matérialiser les modèles en tables.
\end{enumerate}

\subsection*{Transformations attendues}

\begin{itemize}
  \item Typage correct des colonnes.
  \item Conversion des dates.
  \item Conversion des champs numériques avec virgule.
  \item Nettoyage simple des chaînes.
  \item Conservation des colonnes techniques.
  \item Création d’un identifiant technique (clé primaire) pour le personnel
  afin d’éviter l’utilisation directe du \file{staff\_name} comme clé.
\end{itemize}

\subsection*{Tests de qualité}

Mettre en place des tests dbt :

\begin{itemize}
  \item \file{not\_null} sur les identifiants.
  \item \file{unique} sur les clés primaires.
  \item \file{accepted\_values} sur certaines colonnes (ex : accepted, role).
\end{itemize}


% ---------------------------------------------------------------------------------------------

\section{Étape 4 : Chargement incrémental}

\subsection*{Réflexion}

\begin{enumerate}
  \item Quels risques existe-t-il si l’on recharge entièrement les données à chaque exécution ?
  \item Quelle colonne peut servir de watermark ?
  \item Comment gérer la mise à jour d’un enregistrement existant ?
\end{enumerate}

\subsection*{Implémentation}

Adapter le script d’ingestion afin :

\begin{itemize}
  \item de permettre plusieurs batchs,
  \item d’identifier les nouveaux enregistrements,
  \item de conserver l’historique en bronze.
\end{itemize}

Expliquer comment la couche SILVER peut permettre de ne conserver que la version la plus récente.

% ---------------------------------------------------------------------------------------------

\section*{Conclusion}

À l’issue de ce TP, vous devez disposer :

\begin{itemize}
  \item d’une couche BRONZE historisée,
  \item d’une couche SILVER propre et testée,
  \item d’un pipeline ELT reproductible.
\end{itemize}

Ces éléments serviront de base à la construction des indicateurs décisionnels
et du dashboard dans le TP2.

\end{document}